% !TEX root = ../main.tex

% 其他学位学位论文无需此说明

% 学术博士的“创新性说明”
\begin{innovations}
本文围绕粒子物理标准模型的精确检验与新物理探测，基于 ATLAS 实验 Full Run 2 的 $140 \text{ fb}^{-1}$ 质子-质子对撞数据及探测器硬件升级项目，取得了以下具有创新性的研究成果：

\begin{enumerate}
    \item \textbf{实现了迄今最高精度的 $ZZ\rightarrow \ell^+\ell^- \nu \bar\nu$ 产生截面测量。} 通过优化的多区域联合拟合与严格的本底压制策略，将包容性基准截面测量的总不确定度由以往的 $7.0\%$ 显著降低至 $4.6\%$。此外，本研究首次在双轻子加不可见粒子末态中基于标准模型有效场论（SMEFT）框架给出了对中性反常三玻色子耦合（aTGCs）的严格限制，极大拓展了标准模型电弱部分的检验边界。
    \item \textbf{给出了更为先进的暗物质与新物理探测限制。} 在双轻子加不可见末态通道的暗物质搜寻中，引入提升多变量决策树（BDT）算法进行信号分离，在 $95\%$ 置信度（CL）下，将希格斯玻色子不可见衰变分支比 $\mathcal{B}(H \to \text{inv})$ 的观测上限大幅收紧至 $18.7\%$。针对简化暗物质模型及 $\text{2HDM}+a$ 模型给出了更为严苛的排除边界，并在低暗物质质量区间展现出与地下直接探测实验的高度物理互补性。
    \item \textbf{高亮度大型强子对撞机（HL-LHC）缪子探测器升级的关键硬件技术。} 完成了 50 个新型小径监测漂移管（sMDT）高精度探测器的组装与严格性能测试。宇宙射线测试表明其单管空间分辨率优于 $90 \ \mu\text{m}$，完全满足了升级后的HL-LHC的极高本底计数率环境下的苛刻要求，为未来大型强子对撞机的高精度物理运行奠定了坚实的硬件基础。
\end{enumerate}
\end{innovations}


% % 工程博士和专业博士的“创新性及应用性说明”
% \begin{innovations}
%   工程/专业博士学位论文应主要聚焦工程/专业实践和应用研究，
%   须体现创新性、实践性、应用性特征，
%   体现作者在专业领域掌握坚实全面的基础理论和系统深入的专门知识，
%   具有独立承担专业实践工作的能力，在专业实践领域做出创新性成果，
%   对推动本专业领域知识和技术的发展作出重要贡献。
%
%   故在此须说明本学位论文的创新性及应用性，确保符合工程/专业博士学位论文要求，
%   字数要求在 200 至 400 字之间。
% \end{innovations}
